\documentclass[conference]{IEEEtran}
\IEEEoverridecommandlockouts

\usepackage{cite}
\usepackage{amsmath,amssymb,amsfonts}
\usepackage{algorithmic}
\usepackage{algorithm}
\usepackage{graphicx}
\usepackage{textcomp}
\usepackage{xcolor}
\usepackage{booktabs}
\usepackage{array}

\def\BibTeX{{\rm B\kern-.05em{\sc i\kern-.025em b}\kern-.08em
    T\kern-.1667em\lower.7ex\hbox{E}\kern-.125emX}}

\begin{document}

\title{Constant Bandwidth Server:\\Technical Briefing and Scientific Contextualization\\
{\footnotesize Milestones 1 \& 2 --- Real-Time Systems Seminar, HSHL}}

\author{\IEEEauthorblockN{Soaib Ferdous}
\IEEEauthorblockA{\textit{BEng in Electronic Engineering} \\
\textit{Hochschule Hamm-Lippstadt}\\
Student ID: 1232368 \\
Group B, Team 6}}

\maketitle

\begin{abstract}
This document covers Milestones 1 and 2 for the seminar topic
\emph{Constant Bandwidth Server} (CBS).
Milestone 1 provides a structured technical briefing: problem setting,
formal definition, scheduling algorithm, assumptions, runtime overhead,
and open questions.
Milestone 2 places CBS in the landscape of aperiodic server algorithms,
compares it with related approaches (PS, DS, SS, TBS), analyzes
engineering tradeoffs, and discusses consequences for embedded systems.
Primary reference: Buttazzo, \textit{Hard Real-Time Computing Systems}.
\end{abstract}

%%=============================================================
%% MILESTONE 1
%%=============================================================

\section*{--- Milestone 1: Technical Understanding ---}

% ============================================================
\section{Problem Statement and Motivation}
% ============================================================

\begin{itemize}
  \item Real-time systems typically mix \emph{periodic} tasks (hard deadlines,
        known periods) with \emph{aperiodic} tasks (irregular arrivals, often
        soft deadlines).
  \item Under fixed-priority schedulers (RM/DM), naive background execution of
        aperiodic tasks yields poor response times; foreground execution
        can cause periodic tasks to miss deadlines.
  \item Earlier server approaches attempt a middle ground:
        \emph{Polling Server} (PS) wastes budget if no aperiodic job is
        waiting; \emph{Deferrable Server} (DS) preserves budget but can
        violate the RM schedulability bound because it behaves like a task with
        a higher effective utilization at replenishment points.
  \item \textbf{Core problem:} How can aperiodic tasks be served with low
        response time \emph{and} with a provable guarantee that periodic task
        schedulability is not affected?
  \item CBS answers this under EDF by ensuring the server never increases the
        processor demand beyond what is implied by its assigned bandwidth
        $U_s = Q_s / T_s$.
\end{itemize}

% ============================================================
\section{Core Technical Contribution}
% ============================================================

\begin{itemize}
  \item CBS is a dynamic-priority server designed to work with the
        \emph{Earliest Deadline First} (EDF) scheduler.
  \item Key property: \textbf{temporal isolation} --- a misbehaving or
        overloaded aperiodic task cannot steal CPU time from periodic tasks
        beyond what the server's bandwidth allows.
  \item Mechanism: the server has a \emph{current deadline} $d_s$ and a
        \emph{current budget} $c_s$. The aperiodic job runs only when the
        server's deadline $d_s$ is the earliest in the system. Budget is
        consumed during execution. When budget is exhausted, the deadline is
        postponed and budget is recharged --- the task continues but at a
        lower effective priority.
  \item This distinguishes CBS from DS: in DS, the server can retain budget
        across a period boundary, increasing instantaneous demand. CBS
        prevents this by its replenishment rule (Section~\ref{sec:algo}).
\end{itemize}

% ============================================================
\section{Formal Definition}
% ============================================================

A CBS is characterized by two parameters:
\begin{itemize}
  \item $Q_s$ --- maximum server budget (execution budget per period)
  \item $T_s$ --- server period
\end{itemize}

\noindent\textbf{Server bandwidth:}
\begin{equation}
  U_s = \frac{Q_s}{T_s}
  \label{eq:bandwidth}
\end{equation}

\noindent\textbf{Schedulability condition} (CBS + $n$ periodic EDF tasks):
\begin{equation}
  U_s + \sum_{i=1}^{n} \frac{C_i}{T_i} \leq 1
  \label{eq:sched}
\end{equation}
where $C_i$ and $T_i$ are the WCET and period of periodic task $\tau_i$.
CBS is treated as a periodic task with utilization $U_s$ in the EDF
schedulability test, which is its central analytical advantage.

\noindent\textbf{Runtime server state:}
\begin{itemize}
  \item $c_s$: remaining budget, $0 \leq c_s \leq Q_s$
  \item $d_s$: current server deadline
\end{itemize}

% ============================================================
\section{Algorithm and Budget/Deadline Rules}
\label{sec:algo}
% ============================================================

\subsection{Initialization}
At $t = 0$: $c_s \leftarrow Q_s$, $d_s \leftarrow T_s$.

\subsection{Job Arrival Rule}
When aperiodic job $J$ arrives at time $t$:
\begin{itemize}
  \item \textbf{If} server idle and $c_s \geq (d_s - t) \cdot U_s$:
        assign current $d_s$ to $J$ (budget is still ``fresh'').
  \item \textbf{Else:} set $d_s \leftarrow d_s + T_s$,
        $c_s \leftarrow Q_s$; assign new $d_s$ to $J$.
  \item \textbf{If} server busy: $J$ is queued (FIFO).
\end{itemize}

The condition $c_s \geq (d_s - t) \cdot U_s$ prevents using stale budget
that would cause demand to exceed $U_s$ within the current period.

\subsection{Budget Exhaustion Rule}
When $c_s = 0$ at time $t$:
$d_s \leftarrow d_s + T_s$, $c_s \leftarrow Q_s$.
The job continues with the new (later) $d_s$, reducing its EDF priority.

\begin{algorithm}
\caption{CBS decisions (simplified)}
\begin{algorithmic}[1]
\STATE \textbf{On job arrival} at $t$:
\IF{server idle \AND $c_s \geq (d_s - t) \cdot U_s$}
  \STATE assign current $d_s$ to job
\ELSE
  \STATE $d_s \leftarrow d_s + T_s$;\ $c_s \leftarrow Q_s$
  \STATE assign new $d_s$ to job
\ENDIF
\STATE \textbf{On budget exhaustion:}
\STATE $d_s \leftarrow d_s + T_s$;\ $c_s \leftarrow Q_s$
\STATE job continues with updated $d_s$
\end{algorithmic}
\end{algorithm}

% ============================================================
\section{Assumptions and Scope Limitations}
% ============================================================

\begin{itemize}
  \item \textbf{Scheduler:} EDF required. CBS does not apply to
        fixed-priority schedulers (RM/DM).
  \item \textbf{Processor:} Single processor in original formulation.
  \item \textbf{Preemption:} Fully preemptive model.
  \item \textbf{Task model:} Periodic tasks with implicit deadlines
        ($D_i = T_i$), known $C_i$ and $T_i$.
  \item \textbf{Aperiodic guarantees:} Bandwidth only --- no hard
        per-job response time guarantee.
  \item \textbf{No resource sharing:} Priority inversion and shared
        resources are out of scope in basic CBS.
\end{itemize}

% ============================================================
\section{Runtime and Overhead}
% ============================================================

\begin{itemize}
  \item \textbf{Server state update:} $O(1)$ per arrival and per budget
        exhaustion event.
  \item \textbf{EDF queue:} $O(\log n)$ insert/update, shared with all
        tasks.
  \item \textbf{Timer:} One additional timer interrupt per budget period
        when server is active.
  \item \textbf{Memory:} Constant per server: $Q_s$, $T_s$, $c_s$,
        $d_s$, job queue pointer.
\end{itemize}

% ============================================================
\section{Open Questions (Identified at M1)}
% ============================================================

\begin{enumerate}
  \item Hard CBS variant: how does the replenishment rule differ?
  \item Closed-form aperiodic response time bound under CBS?
  \item Exact distinction between CBS and TBS (addressed in M2).
  \item Multiple CBS servers on one processor: schedulability
        generalization?
  \item UPPAAL modeling: how to represent budget and deadline as
        clock constraints in a timed automaton?
\end{enumerate}

%%=============================================================
%% MILESTONE 2
%%=============================================================

\section*{--- Milestone 2: Scientific Contextualization ---}

% ============================================================
\section{Overview of the Aperiodic Server Landscape}
% ============================================================

CBS belongs to a well-defined line of work on aperiodic task servers,
developed primarily by Buttazzo and collaborators over the 1990s--2000s.
The related work is therefore \emph{concentrated}: there is one central
author group, one formal framework (EDF or RM + server), and one line of
incremental improvements. This is not a sparse literature by accident ---
it reflects a mature, self-contained problem class within real-time
scheduling theory.

The servers divide naturally into two families:
\begin{itemize}
  \item \textbf{Fixed-priority family} (work under RM/DM): Polling
        Server (PS), Deferrable Server (DS), Priority Exchange (PX),
        Sporadic Server (SS).
  \item \textbf{Dynamic-priority family} (work under EDF): Total
        Bandwidth Server (TBS), Constant Bandwidth Server (CBS).
\end{itemize}

CBS is the most capable server in the dynamic-priority family and
directly generalizes TBS.

% ============================================================
\section{Comparison of Aperiodic Server Approaches}
% ============================================================

Table~\ref{tab:comparison} summarizes the key properties of the main
aperiodic servers relevant to CBS.

\begin{table*}[t]
\centering
\caption{Comparison of Aperiodic Server Approaches}
\label{tab:comparison}
\renewcommand{\arraystretch}{1.3}
\begin{tabular}{|p{2.2cm}|p{1.8cm}|p{2.0cm}|p{2.5cm}|p{2.5cm}|p{2.8cm}|}
\hline
\textbf{Server} & \textbf{Scheduler} & \textbf{Sched. guarantee} &
\textbf{Overload behavior} & \textbf{Budget rule} &
\textbf{Main limitation} \\
\hline
Polling Server (PS) &
RM (fixed) &
Yes (same as periodic task) &
Budget wasted if no job &
Replenish at period start &
Poor response time when job arrives just after replenishment \\
\hline
Deferrable Server (DS) &
RM (fixed) &
No (can violate RM bound) &
Budget preserved; may overload &
Replenish at period start; budget kept &
Schedulability not provable with standard RM bound \\
\hline
Sporadic Server (SS) &
RM (fixed) &
Yes (behaves like periodic task) &
Budget replenished after execution &
Replenish after $T_s$ from first use &
Complex replenishment tracking; fixed-priority only \\
\hline
Total Bandwidth Server (TBS) &
EDF (dynamic) &
Yes &
Deadline postponed per job &
No explicit budget; assigns deadline $d_k = \max(r_k, d_{k-1}) + C_k/U_s$ &
No isolation: one large job can delay subsequent jobs within server \\
\hline
\textbf{CBS} &
\textbf{EDF (dynamic)} &
\textbf{Yes} &
\textbf{Deadline postponed; budget recharged} &
\textbf{Explicit budget $Q_s$; replenish on exhaustion or stale check} &
\textbf{Requires EDF runtime; single-processor original formulation} \\
\hline
\end{tabular}
\end{table*}

% ============================================================
\section{Detailed Comparison: CBS vs. Neighboring Approaches}
% ============================================================

\subsection{CBS vs. Polling Server}
\begin{itemize}
  \item PS operates under RM. It replenishes budget only at the start
        of each server period. If no aperiodic job is waiting, the budget
        is lost.
  \item \textbf{Tradeoff:} PS has simple, predictable behavior and a
        clean RM schedulability proof, but aperiodic response time can be
        up to one full server period in the worst case.
  \item CBS under EDF achieves much lower response time because the
        server can use available slack immediately, not only at period
        boundaries.
  \item \textbf{Why PS is excluded from deep CBS analysis:} PS is a
        fixed-priority server; the two live in different scheduling
        frameworks. Comparing them requires choosing a common metric
        (e.g., average aperiodic response time at equal utilization),
        which is addressed in Buttazzo's simulation comparisons.
\end{itemize}

\subsection{CBS vs. Deferrable Server}
\begin{itemize}
  \item DS preserves unused budget across period boundaries, giving
        better response time than PS. However, this means DS can
        execute at the start of a new period with a full budget
        \emph{and} has just replenished --- effectively doubling its
        instantaneous demand, which violates the standard RM
        utilization bound.
  \item \textbf{Key tradeoff:} DS improves aperiodic response time
        at the cost of formal schedulability. CBS (and SS) solve the
        schedulability problem DS introduces.
  \item CBS is in a different scheduler family (EDF), so direct
        substitution is not always possible.
\end{itemize}

\subsection{CBS vs. Sporadic Server}
\begin{itemize}
  \item SS solves the same problem as CBS but for \emph{fixed-priority}
        (RM) schedulers. SS replenishes budget $T_s$ time units after
        the first execution unit is consumed, not at a fixed period
        boundary. This gives it RM-compatible schedulability.
  \item \textbf{Tradeoff:} SS requires tracking the replenishment time
        for each budget chunk used, which is more complex to implement
        than CBS. CBS has a simpler replenishment rule (refill to $Q_s$
        on exhaustion or stale arrival) because EDF handles the priority
        ordering automatically.
  \item \textbf{Embedded systems consequence:} If the system already
        uses RM (common in legacy RTOS), SS is the correct CBS analog.
        CBS requires switching to EDF, which may not be feasible on
        existing platforms.
\end{itemize}

\subsection{CBS vs. Total Bandwidth Server}
\begin{itemize}
  \item TBS is the direct EDF predecessor to CBS. TBS assigns each
        arriving aperiodic job $k$ the deadline:
        \begin{equation}
          d_k = \max(r_k,\, d_{k-1}) + \frac{C_k}{U_s}
          \label{eq:tbs}
        \end{equation}
        where $r_k$ is the arrival time and $C_k$ is the actual
        execution time of job $k$.
  \item TBS has no explicit budget variable. It is simpler to implement
        than CBS.
  \item \textbf{Key difference:} TBS uses the \emph{actual} execution
        time $C_k$ of each job, which is only known at job completion.
        This means TBS cannot provide temporal isolation against
        overrunning tasks --- a job that runs longer than expected
        pushes all subsequent deadlines further out. CBS uses an
        \emph{explicit budget} $Q_s$: when the budget runs out, the
        server recharged independently of what the job actually consumed,
        bounding the damage of a single overrunning job.
  \item \textbf{Tradeoff summary:} TBS is simpler and slightly more
        responsive under normal load; CBS provides stronger isolation
        and is safer under overload or in mixed-criticality contexts.
\end{itemize}

% ============================================================
\section{Scientific and Engineering Tradeoffs}
% ============================================================

\begin{itemize}
  \item \textbf{Precision vs. simplicity:} CBS adds an arrival-time
        check ($c_s \geq (d_s - t) \cdot U_s$) absent in TBS. This
        makes CBS more complex but prevents stale-budget overload. For
        simple systems with well-behaved aperiodic tasks, TBS may
        suffice.
  \item \textbf{Formal guarantee vs. practical deployability:} CBS
        requires EDF. Many industrial RTOS (VxWorks, FreeRTOS, classic
        AUTOSAR) use RM/fixed-priority by default. Deploying CBS
        requires either an EDF-capable OS or a custom implementation.
  \item \textbf{Isolation strength vs. utilization efficiency:} CBS's
        budget cap means a short aperiodic burst that fits within $Q_s$
        may still trigger a full $T_s$ wait before a new budget is
        available. TBS in the same scenario would immediately assign a
        tighter deadline. CBS trades some efficiency for stronger
        isolation.
  \item \textbf{Single-core vs. multi-core:} CBS is defined for a
        single processor. Multi-core extension requires partitioning or
        global EDF, both of which introduce new complexity (migration
        overhead, cache effects, scheduling anomalies).
  \item \textbf{Offline vs. online parameter selection:} $Q_s$ and
        $T_s$ must be chosen offline. If aperiodic load is highly
        variable, a fixed $(Q_s, T_s)$ pair may be too conservative
        (wastes bandwidth) or too aggressive (poor isolation). Adaptive
        CBS variants exist but are outside the basic formulation.
\end{itemize}

% ============================================================
\section{Embedded Systems Consequences}
% ============================================================

\begin{itemize}
  \item \textbf{Timing determinism:} CBS provides a formal bound on
        the CPU share consumed by aperiodic tasks. This is essential in
        safety-critical embedded systems (automotive, avionics) where
        interference between task classes must be bounded.
  \item \textbf{Memory:} CBS server state is $O(1)$ per instance.
        Suitable for microcontrollers with limited RAM.
  \item \textbf{Energy:} EDF is work-conserving, which helps with
        dynamic voltage/frequency scaling (DVFS). CBS preserves this
        property since aperiodic tasks run whenever the server has
        budget and processor is free.
  \item \textbf{OS support requirement:} CBS needs an EDF-capable
        scheduler with per-task deadline tracking and a budget timer.
        Standard bare-metal RTOS with only fixed-priority support cannot
        run CBS without modification. POSIX SCHED\_DEADLINE in Linux
        implements a CBS-like mechanism.
  \item \textbf{Deployment example (suitable):} A robotic controller
        with periodic sensor fusion tasks (hard RT) plus event-driven
        network message handlers (soft RT). CBS isolates the network
        handler budget from the control loop.
  \item \textbf{Deployment example (unsuitable):} A legacy automotive
        ECU running AUTOSAR OS with fixed-priority scheduling. CBS
        cannot be applied without replacing the scheduler; SS would be
        the appropriate analog.
\end{itemize}

% ============================================================
\section{Why Some Related Work Is Excluded}
% ============================================================

\begin{itemize}
  \item \textbf{Mixed-criticality scheduling:} Relevant to CBS
        conceptually (task isolation, overload) but addresses a
        different problem (varying assurance levels, WCET uncertainty).
        Not included as related work for CBS specifically.
  \item \textbf{Hierarchical/compositional scheduling:} CBS can be
        used as a component in hierarchical scheduling (each subsystem
        gets a CBS), but the hierarchical framework itself (e.g.,
        Lipari-Bini) is a separate problem class. Relevant to mention
        as an extension, not as direct related work.
  \item \textbf{Non-preemptive scheduling:} CBS assumes full
        preemption. Non-preemptive variants exist but are a different
        design point.
  \item \textbf{Multiprocessor global scheduling papers:} Out of scope
        for the single-processor CBS formulation; relevant only if
        multi-core CBS extension is addressed.
\end{itemize}

%%=============================================================
%% REFERENCES
%%=============================================================

\section*{References}

\begin{thebibliography}{9}

\bibitem{buttazzo}
G.~C.~Buttazzo,
\textit{Hard Real-Time Computing Systems: Predictable Scheduling Algorithms
and Applications}, 4th~ed.
Springer, 2024.
[Primary reference for CBS formal definition, budget and deadline
rules, and schedulability analysis under EDF.]

\bibitem{abeni1998}
L.~Abeni and G.~Buttazzo,
``Integrating Multimedia Applications in Hard Real-Time Systems,''
in \textit{Proc. 19th IEEE Real-Time Systems Symposium
(Cat. No. 98CB36279)}, IEEE, 1998, pp.~4--13.
[Original CBS paper. Introduces the formal definition, budget/deadline
rules, and schedulability proof.]

\bibitem{abeni2004}
L.~Abeni and G.~Buttazzo,
``Resource Reservation in Dynamic Real-Time Systems,''
\textit{Real-Time Systems}, vol.~27, no.~2, pp.~123--167, 2004.
[Extended analysis of CBS including response time results and
comparison with TBS.]

\bibitem{ai_usage_m1}
S.~Ferdous,
\textit{AI Usage Protocol, Milestone 1: Constant Bandwidth Server},
HSHL, 2026.
[Protocol version 0.2; tool: ChatGPT GPT-5.5; 6 entries; JSON and
.bib submitted with milestone.]

\bibitem{ai_usage_m2}
S.~Ferdous,
\textit{AI Usage Protocol, Milestone 2: Constant Bandwidth Server},
HSHL, 2026.
[Protocol version 0.2; tool: Claude claude-sonnet-4-6; 3 entries; JSON and
.bib submitted with milestone.]

\end{thebibliography}

\end{document}
